% !TEX root = ../../tree_notes.tex
% =====================================================================
%  2.1  The Binary Tree Node & Structural Recursion        sec:c02-model
%  Exports: def:c02-treenode  def:c02-recursive-view
%           thm:c02-null-links
%           lst:c02-kit  lst:c02-build
%           prob:c02-second-min  prob:c02-paths  prob:c02-construct
% =====================================================================
\section{The Binary Tree Node and Structural Recursion}\label{sec:c02-model}

\begin{topicintro}
Every binary-tree interview answer is one recursive function over a \code{TreeNode}. This
section isolates the bare skeleton that all of them share -- handle the \nil{} base case,
recurse on both children, combine the two answers with the current node -- and proves it
correct by induction. Once the skeleton is automatic, the rest of the chapter is only about
\emph{what} you combine.

\textbf{Prerequisites.} \secref{c01-what-is-tree} (parent, child, subtree, \nil) and
\secref{c01-anatomy} (the height and size recurrences, which are our first examples).

\textbf{After this section you can}
\begin{itemize}[nosep]
  \item represent a binary tree with a \code{TreeNode} struct and link nodes by hand;
  \item state and prove the "$n$ nodes, $n+1$ null links" counting fact;
  \item write the structural-recursion template and justify it by induction on $\size$;
  \item solve count / collect / compare problems (second-min, root-to-leaf paths,
        parenthesised build) as one template with different combines.
\end{itemize}
\end{topicintro}

\subsection{Intuition}

\begin{analogybox}[A company org chart seen from one manager]
Ask any manager "how many people report to you, directly or indirectly?" They do not walk
the whole company: they ask their two direct reports, get each one's number, add $1$ for
themselves, and reply. That single sentence -- \emph{trust your reports to answer about
their own teams, then add yourself} -- is exactly structural recursion on a tree.
\end{analogybox}

A binary tree node holds a value and points to at most two smaller binary trees. That
self-similarity is the whole game: any question about the tree at a node can be answered
from the \emph{same} question about its left and right subtrees, plus a tiny amount of
work at the node itself.

\begin{figure}[H]
\centering
\begin{tikzpicture}[diagram]
  \node[troot] (r) at (0,0) {$1$};
  \node[tnode] (l) at (-1.7,-1.4) {$2$};
  \node[tnode] (c) at (1.7,-1.4) {$3$};
  \node[tleaf] (ll) at (-2.5,-2.8) {$4$};
  \node[tleaf] (lr) at (-0.9,-2.8) {$5$};
  \node[nullnode] (rr) at (1.7,-2.8) {};
  \draw[tedge] (r) -- (l); \draw[tedge] (r) -- (c);
  \draw[tedge] (l) -- (ll); \draw[tedge] (l) -- (lr);
  \draw[guide] (c) -- (rr);
  \node[annot, text=warningred] at (2.9,-2.8) {null right};
\end{tikzpicture}
\caption{Node $1$ has two children; node $3$ has a left child missing (the dashed guide
line ends at a small open null node). A \code{nullptr} is a real thing you must test for,
not an afterthought.}\label{fig:c02-node}
\end{figure}

\subsection{Definitions and Notation}

\begin{definition}{The \code{TreeNode} model}{c02-treenode}
A binary tree is passed around as a pointer to its root. Each node is a
\code{struct TreeNode \{ int val; TreeNode *left, *right; \};}
where \code{left} and \code{right} are pointers that are \nil{} (\code{nullptr}) when that
child is absent. The tree is the pointer \code{root}; the empty tree is \code{nullptr}. Two
children are \emph{distinguished}: a single left child and a single right child are
different trees (\defref{c01-binary}).
\end{definition}

\begin{definition}{Recursive view of a binary tree}{c02-recursive-view}
A binary tree $T$ is either empty ($\nil$) or a triple $(\text{val}, L, R)$ where $L$ and
$R$ are themselves binary trees (the left and right subtrees). This is the
\emph{inductive} definition: it says a tree is a node glued to two smaller trees, and it is
why every property of $T$ is proved by reasoning about $L$ and $R$ and the node.
\end{definition}

\begin{figure}[H]
\centering
\begin{forest} bintree [1 [2 [4] [5]] [3 [,phantom]]] \end{forest}
\caption{The same tree as \figref{c02-node}, drawn by \texttt{forest} from its recursive
specification \code{[1 [2 [4] [5]] [3 [,phantom]]]}: read a bracket as "a node followed by
its subtrees". The \code{[,phantom]} reserves the missing left slot of $3$.}\label{fig:c02-recursive}
\end{figure}

\subsection{How Many Empty Links?}

Before algorithms, one counting fact that explains why \nil{} handling is unavoidable.

\begin{theorem}{Null Links in a Binary Tree}{c02-null-links}
A binary tree with $n$ nodes has exactly $n-1$ non-\nil{} child pointers and $n+1$
\nil{} (\code{null}) child pointers.
\end{theorem}
\begin{proofbox}
Each of the $n$ nodes owns two child-pointer slots, so there are $2n$ slots in total.
Every node except the root is pointed to by exactly one non-\nil{} child pointer (a tree
has one parent per non-root node and, by \thmref{c01-unique-path}, no node is pointed at
twice), so $n-1$ slots are filled. The remaining slots are null:
$2n-(n-1)=n+1$.
\end{proofbox}
\begin{insightbox}
There is always at least one more null pointer than there are nodes. That is precisely why
"what does the code do when the child is \nil?" is the number-one source of wrong answers:
your recursion meets $n+1$ nulls but only $n$ real nodes, and each null is a base case you
must define. Threaded and array representations (\chref{c04}, \chref{c07}) exist largely to
recycle these $n+1$ dead pointers.
\end{insightbox}

\vsbox{Leaf, null, and "missing one child"}{A leaf}{A null child}{a real node \code{r} with
\code{r->left == r->right == nullptr}; it has a value}{the absent pointer \code{nullptr}; it
has no value and is not a node}

\subsection{The Structural-Recursion Template}

Strip away the problem and almost every solution has this skeleton: act on the base case,
recurse on both subtrees, combine.

\begin{lstlisting}[caption={The structural-recursion kit: count nodes, count leaves, and sum values, all with the same skeleton}, label={lst:c02-kit}]
#include <bits/stdc++.h>
using namespace std;

struct TreeNode {
    int val;
    TreeNode *left, *right;
    TreeNode(int x = 0) : val(x), left(nullptr), right(nullptr) {}
};

// "Trust the recursion": each function only decides the node, given correct answers
// for the two subtrees. The null base case is what the n+1 nulls force us to define.
int countNodes(TreeNode* r) {                 // size: number of nodes
    if (r == nullptr) return 0;               // base case: empty tree contributes 0
    return 1 + countNodes(r->left) + countNodes(r->right);
}

int countLeaves(TreeNode* r) {                // leaves: nodes with two null children
    if (r == nullptr) return 0;
    if (r->left == nullptr && r->right == nullptr) return 1;   // r itself is a leaf
    return countLeaves(r->left) + countLeaves(r->right);
}

long long sumValues(TreeNode* r) {            // fold val over the whole tree
    if (r == nullptr) return 0;
    return r->val + sumValues(r->left) + sumValues(r->right);
}
\end{lstlisting}

\begin{lstlisting}[caption={Building a small tree by linking pointers -- the tree of Figures 2.1 and 2.2}, label={lst:c02-build}]
TreeNode* buildExample() {
    TreeNode* root = new TreeNode(1);
    root->left  = new TreeNode(2);     // 2 becomes left child of 1
    root->right = new TreeNode(3);     // 3 becomes right child of 1
    root->left->left  = new TreeNode(4);
    root->left->right = new TreeNode(5);
    // node 3 keeps left = right = nullptr automatically (its null links)
    return root;
}
\end{lstlisting}

\begin{complexitybox}
All three functions are $\bigO{n}$ time: the recursion tree of the traversal has one call
per node plus the $n+1$ null base cases, so $2n+1$ calls total. Memory is $\bigO{h}$ for the
call stack -- $\bigO{\log n}$ when balanced but $\bigO{n}$ when skewed; a path of $10^5$
nodes is the stack-overflow warning of \secref{c01-anatomy} again.
\end{complexitybox}

\subsection{Worked Examples}

\begin{example}{All three counts on one tree}{c02-kit-ex}
For the tree of \figref{c02-recursive} ($n=5$): \code{countNodes} $=5$;
\code{countLeaves} $=3$ (nodes $4,5$ and $3$, since $3$'s children are both null);
\code{sumValues} $=1+2+3+4+5=15$. The null-link lemma predicts $n+1=6$ nulls; check the
picture: $4,5,3$ each contribute two nulls $=6$.
\end{example}

\begin{tracebox}
\code{countNodes} on a two-node left chain root$\to$a returns
$1+\text{countNodes}(a)+0 = 1+(1+0+0)=2$: the leaf $a$ answers $1$ first, then the root
adds itself -- the aggregate rising from the leaves, the same flow as \thmref{c01-height-recursion}.
\end{tracebox}

\subsection{Common Pitfalls}

\begin{warningbox}[Testing the child, not the pointer]
Guard with \code{if (r == nullptr) return ...;} \emph{before} touching \code{r->val}.
Reading \code{r->left} on a null \code{r} is undefined behaviour (a crash), not a false
condition. The base case comes first, always.
\end{warningbox}

\begin{warningbox}[Confusing "leaf" with "null"]
A leaf is a node; a null is the absence of one. \code{countLeaves} returns $1$ only when
\emph{both children are null}, and must return $0$ when \emph{the node itself is null}.
Swapping these two base cases is the classic one-line bug.
\end{warningbox}

\subsection{Solved Problems}

\solvedproblem{c02-second-min}{Second Minimum Node In a Binary Tree}{LeetCode 671}
  {https://leetcode.com/problems/second-minimum-node-in-a-binary-tree/}{LC Easy}
  {structural recursion, min-combine, sentinel}

\begin{problemstatement}[Second Minimum Node]
Every node here satisfies $\text{val} = \min$ of its two children, so the root holds the
global minimum. Find the second-smallest \emph{distinct} value in the tree, or return $-1$
if none exists.
\end{problemstatement}

\solstep{Understanding the Problem}
The root is the smallest value; the second minimum is the smallest value strictly larger
than $\text{root}\to\text{val}$. Because a subtree rooted at $r$ only ever holds values
$\ge \text{root}\to\text{val}$, we can just scan every node and keep the minimum value that
exceeds the root's -- a one-value combine over the recursion of \coderef{c02-kit}.

\solstep{Key Observations}
\begin{observation}
Any node with $\text{val} > \text{root}\to\text{val}$ is a candidate for the second minimum,
and no descendant can be smaller than it, so we can stop descending there.
\end{observation}
\begin{observation}
The answer is the smallest candidate seen, or $-1$ if no node ever exceeds the root.
\end{observation}

\solstep{Solution}
Carry the fixed root value; recurse, updating a running best over nodes strictly greater
than it and pruning below the first such node.

\solstep{Implementation}
\begin{lstlisting}[caption={Second Minimum Node (LeetCode 671)}, label={lst:c02-second-min}]
class Solution {
    long long best;
    void dfs(TreeNode* r, int rootVal) {
        if (!r) return;
        if (r->val > rootVal) { best = min(best, (long long)r->val); return; }
        dfs(r->left, rootVal);      // r->val == rootVal, keep looking below
        dfs(r->right, rootVal);
    }
public:
    int findSecondMinimumValue(TreeNode* root) {
        best = LLONG_MAX;
        dfs(root, root->val);
        return best == LLONG_MAX ? -1 : (int)best;
    }
};
\end{lstlisting}

\solstep{Complexity}
\begin{complexitybox}
Time $\bigO{n}$ (pruning makes it faster on typical inputs; worst case all-equal visits
every node). Memory $\bigO{h}$ recursion stack. Using \code{long long}/\code{LLONG\_MAX}
for the sentinel avoids colliding with a legitimate value.
\end{complexitybox}

\solstep{Pitfalls and Variants}
\begin{warningbox}[Sentinel that can equal a real value]
Initialising the answer with $0$ or $\texttt{INT\_MAX}$ breaks when those are valid node
values. Use \code{LLONG\_MAX} and translate to $-1$ only at the end.
\end{warningbox}

\solvedproblem{c02-paths}{Binary Tree Paths}{LeetCode 257}
  {https://leetcode.com/problems/binary-tree-paths/}{LC Easy}
  {structural recursion, string build, collect at leaves}

\begin{problemstatement}[Binary Tree Paths]
Return every root-to-leaf path as a string of the form \code{"1->2->3"}. A leaf is a node
with no children; every root-to-leaf path must appear.
\end{problemstatement}

\solstep{Understanding the Problem}
This is the recursive view directly: the paths through a node are its value glued onto each
path of the left and right subtrees, and a leaf contributes exactly one finished path. The
combine is string concatenation, and we collect only at leaves -- the same "descend, then
assemble" shape as \coderef{c02-kit}, just returning a vector.

\solstep{Key Observations}
\begin{observation}
A node's answer $=$ its value prepended to (to) every path from its children.
\end{observation}
\begin{observation}
A leaf's answer is the single string of its own value; a null's answer is the empty list.
\end{observation}

\solstep{Implementation}
\begin{lstlisting}[caption={Binary Tree Paths (LeetCode 257)}, label={lst:c02-paths}]
class Solution {
public:
    vector<string> binaryTreePaths(TreeNode* r) {
        vector<string> out;
        if (!r) return out;
        if (!r->left && !r->right) { out.push_back(to_string(r->val)); return out; }
        for (auto& s : binaryTreePaths(r->left))  out.push_back(to_string(r->val) + "->" + s);
        for (auto& s : binaryTreePaths(r->right)) out.push_back(to_string(r->val) + "->" + s);
        return out;
    }
};
\end{lstlisting}

\solstep{Complexity}
\begin{complexitybox}
Time $\bigO{n \cdot h}$ in the worst case (each of the $\bigO{n}$ path strings can be
$\bigO{h}$ long); the accumulate-a-string-in-one-DFS variant with backtracking is also
$\bigO{n}$ nodes visited. Memory $\bigO{n \cdot h}$ to store all paths. Passing the
running string by value (as above) copies it -- a reference plus explicit backtracking
saves the copies on deep trees.
\end{complexitybox}

\solvedproblem{c02-construct}{Construct String from Binary Tree}{LeetCode 606}
  {https://leetcode.com/problems/construct-string-from-binary-tree/}{LC Easy}
  {structural recursion, parentheses, empty-left rule}

\begin{problemstatement}[Construct String]
Encode a tree as \code{value(left)(right)} with two rules: a leaf prints just its value; if
the left child is null but the right is not, print \code{() } for the missing left, e.g.
\code{"1(2)(())"} style is disallowed -- you must keep the empty \code{()} when only a right
child exists.
\end{problemstatement}

\solstep{Understanding the Problem}
Another pure combine: \code{recurse left} and \code{recurse right} and splice the strings in
with parentheses. The only subtlety is the empty-left exception -- it exists precisely
because without it a "left-only" and a "right-only" one-child node would encode identically,
destroying the distinguished-children rule of \defref{c01-binary}.

\solstep{Implementation}
\begin{lstlisting}[caption={Construct String from Binary Tree (LeetCode 606)}, label={lst:c02-construct}]
class Solution {
public:
    string tree2str(TreeNode* r) {
        if (!r) return "";
        string s = to_string(r->val);
        if (r->left || r->right) s += "(" + tree2str(r->left) + ")";     // left always shown when a child exists
        if (r->right)            s += "(" + tree2str(r->right) + ")";    // right only when present
        return s;
    }
};
\end{lstlisting}

\solstep{Complexity}
\begin{complexitybox}
Time $\bigO{n^2}$ worst case because each \code{+} copies a growing string on a skewed tree;
$\bigO{n}$ if you append into a single \code{string} buffer passed by reference. Memory
$\bigO{h}$ stack plus the output. String-building problems are exactly where
append-in-place beats concatenate-and-return.
\end{complexitybox}

\subsection{Practice Problems}

\begin{practiceset}
  \practice{Maximum Difference Between Node and Ancestor}{LeetCode 1026}
    {https://leetcode.com/problems/maximum-difference-between-node-and-ancestor/}{LC Medium}
    {return the min and max value in each subtree bottom-up; update a global answer with the spread at every node.}
  \practice{Most Frequent Subtree Sum}{LeetCode 508}
    {https://leetcode.com/problems/most-frequent-subtree-sum/}{LC Medium}
    {post-order: a node's subtree sum is val + left-sum + right-sum; tally sums in a hash map.}
  \practicehh{Lowest Common Ancestor of Deepest Leaves}{LeetCode 1123}
    {https://leetcode.com/problems/lowest-common-ancestor-of-deepest-leaves/}{LC Medium}
    {compute each subtree height bottom-up; the LCA is where the two deepest sides tie.}
    {return (height, lca) pairs; if left height > right, the lca bubbles up from the left, and so on.}
\end{practiceset}
